\subsection{Tensor Products and Universal Properties}

% Role: bilinear maps, the tensor-product universal property, its quotient
% construction, uniqueness, and the bilinear-linear correspondence.

Tensor products linearize bilinear data.\\
Their construction combines free extension from Theorem \ref{thm:free-module-map-extension} with quotient descent from Theorem \ref{thm:quotient-universal-property}.

\begin{definition} \label{def:bilinear-map} \hlt{Bilinear Maps}\\
Let $\mathcal{U},\mathcal{V},\mathcal{W}$ be $\mathbb{F}$-vector spaces and $f:\mathcal{U}\times\mathcal{V}\to\mathcal{W}$.
\begin{enumerate}[label=(\roman*)]
\setlength{\itemsep}{0pt}
\item The map $f$ is \hlt{bilinear} if $(\forall\,\mathbf{v}\in\mathcal{V})\,f(\,\cdot\,,\mathbf{v})\in\LL(\mathcal{U},\mathcal{W})$ and $(\forall\,\mathbf{u}\in\mathcal{U})\,f(\mathbf{u},\,\cdot\,)\in\LL(\mathcal{V},\mathcal{W})$.
\item $\operatorname{Bil}_{\mathbb{F}}(\mathcal{U},\mathcal{V};\mathcal{W})\defeq\{f:\mathcal{U}\times\mathcal{V}\to\mathcal{W}:f\text{ is bilinear}\}$.
\item If $\mathcal{W}=\mathbb{F}$, then $\operatorname{Bil}_{\mathbb{F}}(\mathcal{U},\mathcal{V};\mathbb{F})$ is the set of \hlt{bilinear forms} from Definition \ref{def:bilinear-form-and-matrix}.
\end{enumerate}
\end{definition}

\begin{example} \label{ex:basic-bilinear-maps} \hlt{Basic Bilinear Maps}
\begin{enumerate}[label=(\roman*)]
\setlength{\itemsep}{0pt}
\item For an $\mathbb{F}$-algebra $A$, define $\mu:A\times A\to A$ by $\mu(a,b)\defeq ab$.\\
The map $\mu$ is bilinear by distributivity and compatibility of scalar and algebra multiplication.
\item Define $\operatorname{ev}:\LL(\mathcal{V},\mathcal{W})\times\mathcal{V}\to\mathcal{W}$ by $\operatorname{ev}(T,\mathbf{v})\defeq T\mathbf{v}$. The map $\operatorname{ev}$ is bilinear by pointwise linearity in $T$ and the linearity of $T$ in $\mathbf{v}$.
\item For $\varphi\in\mathcal{V}'$, $\psi\in\mathcal{W}'$, define $B_{\varphi,\psi}:\mathcal{V}\times\mathcal{W}\to\mathbb{F}$ by $B_{\varphi,\psi}(\mathbf{v},\mathbf{w})\defeq\varphi(\mathbf{v})\psi(\mathbf{w})$. Map $B_{\varphi,\psi}$ is bilinear by linearity of $\varphi$, $\psi$.
\end{enumerate}
\end{example}

\begin{definition} \label{def:tensor-product-universal-property} \hlt{Tensor Product by Universal Property}\\
Let $\mathcal{U},\mathcal{V},\mathcal{T}$ be $\mathbb{F}$-vector spaces and $t\in\operatorname{Bil}_{\mathbb{F}}(\mathcal{U},\mathcal{V};\mathcal{T})$.\\
Given any $\mathbb{F}$-vector space $\mathcal{W}$, define $t_{\mathcal{W}}^\ast:\LL(\mathcal{T},\mathcal{W})\to\operatorname{Bil}_{\mathbb{F}}(\mathcal{U},\mathcal{V};\mathcal{W})$ by $t_{\mathcal{W}}^\ast(L)\defeq L\circ t$.
\begin{enumerate}[label=(\roman*)]
\setlength{\itemsep}{0pt}
\item The pair $(\mathcal{T},t)$ is a \hlt{tensor product} of $\mathcal{U}$ and $\mathcal{V}$ if every map $t_{\mathcal{W}}^\ast$ above is bijective.
\item For $f\in\operatorname{Bil}_{\mathbb{F}}(\mathcal{U},\mathcal{V};\mathcal{W})$, the map $\widetilde{f}\defeq(t_{\mathcal{W}}^\ast)^{-1}(f)$ is the \hlt{mediating map} for $f$; equivalently, $\widetilde{f}\circ t=f$.
\item For a chosen tensor product, write $\mathcal{T}=\mathcal{U}\otimes\mathcal{V}$ and $\mathbf{u}\otimes\mathbf{v}\defeq t(\mathbf{u},\mathbf{v})$. A tensor of the form $\mathbf{u}\otimes\mathbf{v}$ is \hlt{decomposable}.
\end{enumerate}
\end{definition}

\begin{illustrationfigure}[H]
  \centering
  % Source note: Roman, Advanced Linear Algebra, Ch. 14, Fig. 14.4.
% Redrawn as a schematic universal-property diagram for these notes.
\begin{tikzpicture}[
  >=Stealth,
  every node/.style={font=\small},
  arrow/.style={->, line width=0.55pt},
  factor/.style={->, dashed, line width=0.55pt},
  note/.style={font=\scriptsize, text=gray!70!black}
]
  \node (uv) at (0,1.25) {$\mathcal{U}\times\mathcal{V}$};
  \node (t) at (3.6,1.25) {$\mathcal{U}\otimes\mathcal{V}$};
  \node (w) at (3.6,-0.8) {$\mathcal{W}$};

  \draw[arrow] (uv) -- node[above, note] {$t$ bilinear} (t);
  \draw[arrow] (uv) -- node[below left, note] {$f$ bilinear} (w);
  \draw[factor] (t) -- node[right, note] {$\widetilde{f}$ linear} (w);
\end{tikzpicture}

  \caption{A tensor product is universal for bilinear maps out of $\mathcal{U}\times\mathcal{V}$.}
  \label{fig:tensor-universal-property}
\end{illustrationfigure}

Figure \ref{fig:tensor-universal-property} depicts $f=\widetilde{f}\circ t$ from Definition \ref{def:tensor-product-universal-property}. The definition owns existence and uniqueness of $\widetilde{f}$.

\begin{theorem} \label{thm:tensor-product-quotient-construction} \hlt{Quotient Construction of the Tensor Product}\\
Let $\mathcal{U},\mathcal{V}$ be $\mathbb{F}$-vector spaces, $\mathcal{F}\defeq\bigoplus\limits_{(\mathbf{u},\mathbf{v})\in\mathcal{U}\times\mathcal{V}}^{\mathrm{ext}}\mathbb{F}$. For $(\mathbf{u},\mathbf{v})\in\mathcal{U}\times\mathcal{V}$, let $[\mathbf{u},\mathbf{v}]\in\mathcal{F}$ be coordinate family supported at $(\mathbf{u},\mathbf{v})$ with value $1$.\\
In the following relation families, $r,s\in\mathbb{F}$, $\mathbf{u},\mathbf{u}_1,\mathbf{u}_2\in\mathcal{U}$, and $\mathbf{v},\mathbf{v}_1,\mathbf{v}_2\in\mathcal{V}$.\\
Let $\mathcal{R}_{\mathcal{U}}\leq\mathcal{F}$ be spanned by all $r[\mathbf{u}_1,\mathbf{v}]+s[\mathbf{u}_2,\mathbf{v}]-[r\mathbf{u}_1+s\mathbf{u}_2,\mathbf{v}]$. Let $\mathcal{R}_{\mathcal{V}}\leq\mathcal{F}$ be spanned by all $r[\mathbf{u},\mathbf{v}_1]+s[\mathbf{u},\mathbf{v}_2]-[\mathbf{u},r\mathbf{v}_1+s\mathbf{v}_2]$.\\
Set $\mathcal{R}\defeq\mathcal{R}_{\mathcal{U}}+\mathcal{R}_{\mathcal{V}}$, $\mathcal{U}\otimes\mathcal{V}\defeq\mathcal{F}/\mathcal{R}$, and $\mathbf{u}\otimes\mathbf{v}\defeq[\mathbf{u},\mathbf{v}]+\mathcal{R}$. Define $t:\mathcal{U}\times\mathcal{V}\to\mathcal{U}\otimes\mathcal{V}$ by $t(\mathbf{u},\mathbf{v})\defeq\mathbf{u}\otimes\mathbf{v}$. Then:
\begin{enumerate}[label=(\roman*)]
\setlength{\itemsep}{0pt}
\item $\mathcal{U}\otimes\mathcal{V}=\operatorname{span}\{\mathbf{u}\otimes\mathbf{v}:\mathbf{u}\in\mathcal{U},\ \mathbf{v}\in\mathcal{V}\}$.
\item $(\mathcal{U}\otimes\mathcal{V},t)$ is a tensor product of $\mathcal{U}$ and $\mathcal{V}$.
\end{enumerate}
\begin{proof}
$\mathcal{F}$ is an $\mathbb{F}$-vector space by Proposition \ref{prop:external-direct-sum-family-vector-space}(ii). Let $\mathcal{B}\defeq\{[\mathbf{u},\mathbf{v}]:(\mathbf{u},\mathbf{v})\in\mathcal{U}\times\mathcal{V}\}$.\\
Every $\mathbf{x}=(x_{\mathbf{u},\mathbf{v}})\in\mathcal{F}$ has a unique finite expansion $\mathbf{x}=\sum\limits_{(\mathbf{u},\mathbf{v})\in\operatorname{supp}(\mathbf{x})}x_{\mathbf{u},\mathbf{v}}[\mathbf{u},\mathbf{v}]$ by Definition \ref{def:external-products-and-direct-sums-family}.\\
The set $\mathcal{B}$ is a basis by this expansion, and $\mathcal{F}$ is free on $\mathcal{B}$ by Definition \ref{def:free-module-basis}.\\
Let $j:\mathcal{U}\times\mathcal{V}\to\mathcal{F}$ satisfy $j(\mathbf{u},\mathbf{v})\defeq[\mathbf{u},\mathbf{v}]$. Let $\pi:\mathcal{F}\to\mathcal{U}\otimes\mathcal{V}$ be the quotient map.
\begin{enumerate}[label=(\roman*)]
\setlength{\itemsep}{0pt}
\item Set $\mathcal{D}\defeq\{\mathbf{u}\otimes\mathbf{v}:\mathbf{u}\in\mathcal{U},\ \mathbf{v}\in\mathcal{V}\}$. Then $\pi(\mathcal{B})=\mathcal{D}$ by the definitions of $\pi$ and decomposable tensors.\\
The equality $\mathcal{U}\otimes\mathcal{V}=\pi(\mathcal{F})=\operatorname{span}(\pi(\mathcal{B}))=\operatorname{span}(\mathcal{D})$ holds by the linearity and surjectivity of $\pi$.
\item Fix $r,s\in\mathbb{F}$, $\mathbf{u}_1,\mathbf{u}_2\in\mathcal{U}$, and $\mathbf{v}\in\mathcal{V}$. The identity $t(r\mathbf{u}_1+s\mathbf{u}_2,\mathbf{v})=rt(\mathbf{u}_1,\mathbf{v})+st(\mathbf{u}_2,\mathbf{v})$ holds by the definition of $\mathcal{R}_{\mathcal{U}}$.\\
Fix $\mathbf{u}\in\mathcal{U}$ and $\mathbf{v}_1,\mathbf{v}_2\in\mathcal{V}$. The identity $t(\mathbf{u},r\mathbf{v}_1+s\mathbf{v}_2)=rt(\mathbf{u},\mathbf{v}_1)+st(\mathbf{u},\mathbf{v}_2)$ holds by the definition of $\mathcal{R}_{\mathcal{V}}$.\\
The map $t=\pi\circ j$ is bilinear by these identities. Fix an $\mathbb{F}$-vector space $\mathcal{W}$ and $f\in\operatorname{Bil}_{\mathbb{F}}(\mathcal{U},\mathcal{V};\mathcal{W})$.\\
The assignment $[\mathbf{u},\mathbf{v}]\mapsto f(\mathbf{u},\mathbf{v})$ extends uniquely to $\alpha\in\LL(\mathcal{F},\mathcal{W})$ by Theorem \ref{thm:free-module-map-extension}.\\
Fix a displayed generator $x_{\mathcal{U}}\defeq r[\mathbf{u}_1,\mathbf{v}]+s[\mathbf{u}_2,\mathbf{v}]-[r\mathbf{u}_1+s\mathbf{u}_2,\mathbf{v}]$ of $\mathcal{R}_{\mathcal{U}}$.\\
Then $\alpha(x_{\mathcal{U}})=rf(\mathbf{u}_1,\mathbf{v})+sf(\mathbf{u}_2,\mathbf{v})-f(r\mathbf{u}_1+s\mathbf{u}_2,\mathbf{v})=0$ by bilinearity of $f$.\\
Fix a displayed generator $x_{\mathcal{V}}\defeq r[\mathbf{u},\mathbf{v}_1]+s[\mathbf{u},\mathbf{v}_2]-[\mathbf{u},r\mathbf{v}_1+s\mathbf{v}_2]$ of $\mathcal{R}_{\mathcal{V}}$.\\
Then $\alpha(x_{\mathcal{V}})=rf(\mathbf{u},\mathbf{v}_1)+sf(\mathbf{u},\mathbf{v}_2)-f(\mathbf{u},r\mathbf{v}_1+s\mathbf{v}_2)=0$ by bilinearity of $f$.\\
$\mathcal{R}\leq\ker\alpha$ by the linearity of $\alpha$ and the two generator calculations.\\
Unique $\widetilde{f}\in\LL(\mathcal{U}\otimes\mathcal{V},\mathcal{W})$ satisfies $\widetilde{f}\circ\pi=\alpha$ by Theorem \ref{thm:quotient-universal-property}(i). Equality $\widetilde{f}\circ t=f$ holds by $(\widetilde{f}\circ t)(\mathbf{u},\mathbf{v})=\alpha[\mathbf{u},\mathbf{v}]=f(\mathbf{u},\mathbf{v})$.\\
If $L\in\LL(\mathcal{U}\otimes\mathcal{V},\mathcal{W})$ satisfies $L\circ t=f$, then $L=\widetilde{f}$ by agreement on the spanning decomposable tensors from (i).
\end{enumerate}
\hlt{Structural Commentary}: Free extension kills both relation subspaces and descends through their quotient.\\
\hlt{Key Insight}: The two relation subspaces encode linearity in the tensor factors.
\end{proof}
\end{theorem}

\begin{illustrationfigure}[H]
  \centering
  % Source note: Roman, Advanced Linear Algebra, Ch. 14, Fig. 14.5.
% Redrawn to show the free-vector-space quotient construction.
\begin{tikzpicture}[
  >=Stealth,
  x=0.9cm,
  y=0.9cm,
  every node/.style={font=\small},
  arrow/.style={->, line width=0.65pt, black},
  aux/.style={->, dashed, line width=0.65pt, black},
  factor/.style={->, dashed, line width=0.7pt, black},
  note/.style={font=\scriptsize, text=black, fill=white, inner sep=0.8pt}
]
  \node (uv) at (0,0.45) {$\mathcal{U}\times\mathcal{V}$};
  \node (free) at (3.2,0.45) {$\mathbb{F}^{(\mathcal{U}\times\mathcal{V})}$};
  \node (quot) at (6.55,0.45) {$\mathcal{U}\otimes\mathcal{V}$};
  \node (w) at (3.2,-0.58) {$\mathcal{W}$};

  \draw[arrow] (uv) -- node[above, note] {$j$} (free);
  \draw[arrow] (free) -- node[above, note, pos=0.55] {$\pi$} (quot);
  \draw[arrow] (uv.north east) to[out=30,in=150]
    node[above, note] {$t=\pi j$} (quot.north west);
  \draw[arrow] (uv) -- node[below left, note] {$f$} (w);
  \draw[aux] (free) -- node[right, note, xshift=2pt] {$\alpha$} (w);
  \draw[factor] (quot) -- node[right, note, xshift=8pt] {$\widetilde{f}$} (w);
\end{tikzpicture}

  \caption{The free linear map $\alpha$ kills $\mathcal{R}$, so it descends uniquely to $\widetilde{f}:\mathcal{U}\otimes\mathcal{V}\to\mathcal{W}$.}
  \label{fig:tensor-quotient-construction-universal}
\end{illustrationfigure}

Figure \ref{fig:tensor-quotient-construction-universal} records $j(\mathbf{u},\mathbf{v})=[\mathbf{u},\mathbf{v}]$, $t=\pi\circ j$, and $\widetilde{f}\circ\pi=\alpha$ from Theorem \ref{thm:tensor-product-quotient-construction}.\\
The theorem owns the spanning and factorisation claims.

\begin{corollary} \label{cor:tensor-product-unique-up-to-isomorphism} \hlt{Tensor Products Are Unique up to Isomorphism}\\
If $(\mathcal{T},t)$ and $(\mathcal{S},s)$ are tensor products of $\mathcal{U}$ and $\mathcal{V}$, then a unique isomorphism $\Phi:\mathcal{T}\to\mathcal{S}$ satisfies $\Phi\circ t=s$.
\begin{proof}
Unique mediating maps $\Phi:\mathcal{T}\to\mathcal{S}$ and $\Psi:\mathcal{S}\to\mathcal{T}$ exist by Definition \ref{def:tensor-product-universal-property}(i). They satisfy $\Phi\circ t=s$ and $\Psi\circ s=t$.\\
The identity $\Psi\Phi=I_{\mathcal{T}}$ holds by mediator uniqueness because $(\Psi\Phi)\circ t=\Psi\circ s=t=I_{\mathcal{T}}\circ t$.\\
Likewise, $\Phi\Psi=I_{\mathcal{S}}$ by mediator uniqueness because $(\Phi\Psi)\circ s=\Phi\circ t=s=I_{\mathcal{S}}\circ s$.\\
The map $\Phi$ is an isomorphism with inverse $\Psi$ by these two identities.\\
If $\Theta:\mathcal{T}\to\mathcal{S}$ is linear and $\Theta\circ t=s$, then $\Theta=\Phi$ by uniqueness of the mediating map for $s$.\\
\hlt{Structural Commentary}: The two universal factorisations compose to the unique mediating identity maps.\\
\hlt{Key Insight}: Mediator uniqueness forces the comparison maps to be mutual inverses.
\end{proof}
\end{corollary}

\begin{theorem} \label{thm:tensor-linearizes-bilinear-maps} \hlt{Bilinearity Equals Linearity on a Tensor Product}\\
Let $\mathcal{U},\mathcal{V},\mathcal{W}$ be $\mathbb{F}$-vector spaces. Fix a tensor product $(\mathcal{U}\otimes\mathcal{V},t)$ of $\mathcal{U}$ and $\mathcal{V}$.
\begin{enumerate}[label=(\roman*)]
\setlength{\itemsep}{0pt}
\item The set $\operatorname{Bil}_{\mathbb{F}}(\mathcal{U},\mathcal{V};\mathcal{W})$ is an $\mathbb{F}$-vector space under pointwise addition and scalar multiplication.
\item The assignment $f\mapsto\widetilde{f}$ defines an isomorphism $\Phi:\operatorname{Bil}_{\mathbb{F}}(\mathcal{U},\mathcal{V};\mathcal{W})\to\LL(\mathcal{U}\otimes\mathcal{V},\mathcal{W})$.\\
For $f\in\operatorname{Bil}_{\mathbb{F}}(\mathcal{U},\mathcal{V};\mathcal{W})$, the map $\Phi(f)=\widetilde{f}$ is characterized by $\widetilde{f}\circ t=f$.
\end{enumerate}
\begin{proof}
\begin{enumerate}[label=(\roman*)]
\setlength{\itemsep}{0pt}
\item The zero map is bilinear by Definition \ref{def:bilinear-map}(i). Fix $a,b\in\mathbb{F}$ and $f,g\in\operatorname{Bil}_{\mathbb{F}}(\mathcal{U},\mathcal{V};\mathcal{W})$.\\
For $\mathbf{u}\in\mathcal{U}$ and $\mathbf{v}\in\mathcal{V}$, $(af+bg)(\,\cdot\,,\mathbf{v})\in\LL(\mathcal{U},\mathcal{W})$ and $(af+bg)(\mathbf{u},\,\cdot\,)\in\LL(\mathcal{V},\mathcal{W})$ by linearity of the corresponding partial maps.\\
$af+bg\in\operatorname{Bil}_{\mathbb{F}}(\mathcal{U},\mathcal{V};\mathcal{W})$ by Definition \ref{def:bilinear-map}(i). The vector-space axioms hold by their pointwise counterparts in $\mathcal{W}$.
\item The assignment $\Phi(f)=\widetilde{f}$ is well-defined by Definition \ref{def:tensor-product-universal-property}(i). Define $\Psi:\LL(\mathcal{U}\otimes\mathcal{V},\mathcal{W})\to\operatorname{Bil}_{\mathbb{F}}(\mathcal{U},\mathcal{V};\mathcal{W})$ by $\Psi(L)\defeq L\circ t$.\\
The map $\Psi$ is well-defined by linearity of $L$ and bilinearity of $t$.\\
The identity $\Psi(\Phi(f))=f$ holds by $\widetilde{f}\circ t=f$, and $\Phi(\Psi(L))=L$ by mediator uniqueness for $L\circ t$.\\
Hence $\Psi=\Phi^{-1}$ by these identities. Fix $a,b\in\mathbb{F}$ and bilinear $f,g$.\\
The equality $(a\Phi(f)+b\Phi(g))\circ t=af+bg=\Phi(af+bg)\circ t$ holds by the definition of $\Phi$.\\
Thus $\Phi(af+bg)=a\Phi(f)+b\Phi(g)$ by mediator uniqueness. Map $\Phi$ is linear isomorphism by this linearity and inverse above.
\end{enumerate}
\hlt{Structural Commentary}: Mediator uniqueness constructs the inverse and transfers pointwise linear combinations through $t$.\\
\hlt{Key Insight}: Precomposition with $t$ inverts bilinear linearization.
\end{proof}
\end{theorem}

\begin{computationmethod} \label{meth:defining-linear-maps-on-tensor-products} \hlt{Constructing Maps from Tensor Products}\\
Fix $\mathbb{F}$-vector spaces $\mathcal{U},\mathcal{V},\mathcal{W}$ and a tensor product $(\mathcal{U}\otimes\mathcal{V},t)$. To construct $L\in\LL(\mathcal{U}\otimes\mathcal{V},\mathcal{W})$, proceed as follows:
\begin{enumerate}[label=(\roman*)]
\setlength{\itemsep}{0pt}
\item specify a rule $f:\mathcal{U}\times\mathcal{V}\to\mathcal{W}$;
\item verify $(\forall\,\mathbf{v}\in\mathcal{V})\,f(\,\cdot\,,\mathbf{v})\in\LL(\mathcal{U},\mathcal{W})$ and $(\forall\,\mathbf{u}\in\mathcal{U})\,f(\mathbf{u},\,\cdot\,)\in\LL(\mathcal{V},\mathcal{W})$;
\item set $L\defeq\Phi(f)$; then $L\circ t=f$ by Theorem \ref{thm:tensor-linearizes-bilinear-maps}(ii).
\end{enumerate}
The decomposable tensors span $\mathcal{U}\otimes\mathcal{V}$ by Theorem \ref{thm:tensor-product-quotient-construction}(i), but are generally linearly dependent.\\
Assignments on decomposable tensors must therefore respect their linear relations.
\end{computationmethod}
